\section{Penanda Linguistik ETHICS Domain Gender}
\label{sec:lampiran-ethics-gender-lm}

Gambar berikut melengkapi penanda linguistik ETHICS pada Bab VI. Seluruh
hasil masih berstatus \textit{candidate} dan belum menggantikan audit manusia
independen. Koreksi Holm mengikuti keluarga pengujian pada hasil sumber,
bukan dihitung ulang berdasarkan pembagian gambar. Pada grafik langsung,
segitiga menunjukkan $p_{\mathrm{Holm}}<0{,}05$ dan lingkaran menunjukkan
hasil lainnya; warna membedakan model.

\clearpage
\begin{figure}[H]
    \centering
    \includegraphics[width=\textwidth]{image/BabVI/ethics/gender/linguistic_marker_20261002/00_cakupan.png}
    \caption{Ketersediaan penjelasan dan pasangan analisis; L: laki-laki, P: perempuan (\textit{candidate})}
    \label{fig:ethics-gender-lm-cakupan}
\end{figure}

\clearpage
\begin{figure}[H]
    \centering
    \includegraphics[width=\textwidth]{image/BabVI/ethics/gender/linguistic_marker_20261002/02_panjang_penjelasan.png}
    \caption{Panjang penjelasan menurut skenario; titik: median, garis: rentang antarkuartil (\textit{candidate})}
    \label{fig:ethics-gender-lm-panjang}
\end{figure}

\clearpage
\begin{figure}[H]
    \centering
    \includegraphics[width=\textwidth]{image/BabVI/ethics/gender/linguistic_marker_20261002/N01_bahasa_umum_vs_netral.png}
    \caption{Ukuran efek penanda linguistik umum terhadap skenario netral (\textit{candidate})}
    \label{fig:ethics-gender-lm-umum-netral}
\end{figure}

\clearpage
\begin{figure}[H]
    \centering
    \includegraphics[width=\textwidth]{image/BabVI/ethics/gender/linguistic_marker_20261002/N03_wacana_gender_vs_netral.png}
    \caption{Ukuran efek penanda linguistik wacana gender terhadap skenario netral (\textit{candidate})}
    \label{fig:ethics-gender-lm-wacana-netral}
\end{figure}

\clearpage
\begin{figure}[H]
    \centering
    \includegraphics[width=\textwidth]{image/BabVI/ethics/gender/linguistic_marker_20261002/01_bahasa_umum_man_minus_woman.png}
    \caption{Perbedaan penanda linguistik umum antara persona laki-laki dan perempuan (\textit{candidate})}
    \label{fig:ethics-gender-lm-umum-langsung}
\end{figure}

\clearpage
\begin{figure}[H]
    \centering
    \includegraphics[width=\textwidth]{image/BabVI/ethics/gender/linguistic_marker_20261002/03_bahasa_gender_man_minus_woman.png}
    \caption{Perbedaan penanda linguistik wacana gender antara persona laki-laki dan perempuan (\textit{candidate})}
    \label{fig:ethics-gender-lm-wacana-langsung}
\end{figure}

\clearpage
\begin{figure}[H]
    \centering
    \includegraphics[width=\textwidth]{image/BabVI/ethics/gender/linguistic_marker_20261002/04_marker_biner_man_minus_woman.png}
    \caption{Selisih prevalensi penanda biner antara persona laki-laki dan perempuan; garis: interval kepercayaan 95\% marginal (\textit{candidate})}
    \label{fig:ethics-gender-lm-biner-langsung}
\end{figure}

